\scanpage{12}
\mspage{9}

des dérivateurs'' consiste en grande partie, inversement, chaque fois qu'on a un diagramme dans une cat. de coefficients $\mathbb D(I)$ (sur un \deleted{cat.} diagramme $I$), disons diagramme de type $K$ : $K\longrightarrow\mathbb D(I)$, de se demander s'il ne provient pas ``canoniquement'' d'un objet de $\mathbb D(I\times K^{\unsure{\circ|o}})$ i.e. d'un coefficient sur $I\times K^{\circ}$.

Ainsi pour le carré p. 8. On voudrait qu'il provienne de $\mathbb D(I\times J)$ \deleted{\unsure{=}} où $J=\unsure{\mathbb{\Delta}|\Delta}_{1}\times\mathbb{\Delta}_{1}$, \deleted{\unsure{donc}} et ceci \unsure{fonctoriellement|fonctoriellemt}, donc on désire un foncteur ``canonique''

\[
\mathbb D(I)\longrightarrow\mathbb D(I\times\mathbb{\Delta}_{1}\times\mathbb{\Delta}_{1})
\]

qui, en composant avec

\[
\mathbb D(I\times\mathbb{\Delta}_{1}\times\mathbb{\Delta}_{1})^{\deleted{\unsure{\circ}}}\longrightarrow\underline{\mathrm{Hom}}\,((\mathbb{\Delta}_{1}\times\mathbb{\Delta}_{1})^{\circ},\mathbb D(I))
\]

donne le foncteur décrit par le carré (1). De plus, il faut exprimer dans le contexte présent le caractère ``bi-cartésien'' du carré (1), dans le contexte de la page 2. Question analogue pour le carré (\unsure{1|2}) de la page \unsure{2|?} :

\diagram{\begin{tikzcd}[ampersand replacement=\&]
\& \& F\arrow[dr] \& \\ (2) \& i_{*}i^{!}F\arrow[ur]\arrow[dr,very thick] \& \& j_{*}j^{*}F \\ \& \& e\arrow[ur] \& {,}
\end{tikzcd}}

où il faut de plus interpréter le caractère cartésien.

On va définir en fait un \deleted{diagramme}\inserted{coeff.} sur la sous-cat. ouverte $W$ suivante de $\mathbb{\Delta}_{1}\times\mathbb{\Delta}_{1}$

\diagram{\begin{tikzcd}[ampersand replacement=\&]
(1\unsure{;|,}0) \& \\ \& (0,0)\arrow[ul]\arrow[dl,very thick] \\ (0,1) \& {,\ \unsure{a|ou|à}}
\end{tikzcd}}

% archive sheet number: 11
